% Supersedes appendix_window_check.tex (the exogenous-OU window check has
% been dropped in favor of this feedback-coupled check -- it directly
% tests the paper's own novel mechanism rather than an omitted-static-
% field story that the equivalence literature already covers).

\section{Supplementary check: does feedback affect network estimation?}\label{app:network_feedback_check}

The network-estimation check reported in the main text (Section~\ref{sec:res_network_check}) treated the slow context as a fixed, per-person covariate: for a given level of within-group heterogeneity, each simulated person's value of $P$ was drawn independently, with no dependence on that person's own symptom history. This leaves open a more specific question: does the paper's central dynamical mechanism -- sustained symptom elevation feeding back onto the slow context (Eq.~\eqref{eq:slow_update}) -- itself contribute to, or worsen, the network-estimation distortion documented above?

To test this directly, we gave each simulated individual their own private, coupled fast/slow system rather than an exogenously sampled or generated context. At every simulated time step, an individual's full 12-symptom configuration was drawn from the same pairwise model used in the main check, conditional on that person's current $P_t$; their smoothed symptom signal $m_{\text{slow},t}$ was then updated via the same exponentially weighted average used in the main model (Eq.~\eqref{eq:mslow}), and their own $P_t$ was updated via the calibrated slow-state equation (Eq.~\eqref{eq:slow_update}), using the same parameter values as the main scenario simulations ($\kappa=0.20$, $\sigma_P=0.04$, $\lambda_m=0.001$, $m_\star=0.25$). We compared two conditions, otherwise identical: feedback disabled ($b=0$) and feedback enabled at its calibrated value ($b=0.15$). Each individual was simulated for 4{,}000 steps from a common starting point, long enough for the coupled dynamics to reach a stable regime (verified by comparing the mean context and smoothed symptom signal at successive checkpoints within the burn-in); the symptom vector and context value used for network estimation were taken from the final simulated step, exactly as in the main check.

Enabling feedback produced a large, statistically robust upward shift in the high-baseline group's realized context: mean $\bar P$ rose from 0.999 without feedback to 1.36 with it, closely matching the group-level separation documented in the main scenario simulations (Section~\ref{sec:res_B}), where sustained symptom elevation likewise drove the slow state upward in the high-baseline, bistable-capable condition. The low-baseline group shifted only slightly (from $-0.300$ to $-0.340$), because the feedback reference level $m_\star = 0.25$ sits close to this group's own resting symptom-activation rate, leaving little room for feedback to act.

Critically, however, feedback did not meaningfully widen the \emph{within}-group spread of realized context: $SD(P)$ in the high-baseline group changed only from 0.063 to 0.065, a difference of the same order as its own standard error. This is a direct consequence of how heavily the feedback signal is smoothed: $m_{\text{slow},t}$ averages over roughly $1/\lambda_m = 1{,}000$ steps, the same smoothing used throughout the paper to distinguish sustained symptom elevation from momentary fluctuation. That same smoothing means almost all person-to-person sampling variation in momentary symptom activity is averaged away before it can reach $P_t$, so feedback acts as a nearly deterministic, shared pull on every individual in a group rather than a source of additional between-person divergence.

Because network-estimation distortion in this design is driven by within-group heterogeneity in the omitted covariate, not by between-group differences in its mean, this near-absence of added within-group spread predicts that feedback should not appreciably widen the gap between the symptom-only and context-adjusted estimators -- and this is what we found. Estimated global strength for the high-baseline group was 12.5 (symptom-only) versus 12.4 (context-adjusted) without feedback, and 13.6 versus 13.5 with feedback, against a true value of 8.08; the low-baseline group showed the same close correspondence (14.2 vs.\ 14.1 without feedback, 14.5 vs.\ 14.4 with feedback). Recovery error relative to $W^{\mathrm{true}}$ showed the identical pattern (high-baseline group: 0.116 vs.\ 0.116 without feedback, 0.142 vs.\ 0.142 with feedback; Figure~\ref{fig:network_feedback_check}). The modest increase in both estimators under feedback is consistent with the higher, feedback-elevated baseline rate of symptom activation in that condition rather than with an omitted-variable effect specific to the symptom-only estimator, since both estimators move together.

These results indicate that the paper's central dynamical mechanism -- feedback from sustained symptom elevation onto the slow context -- reproduces and reinforces the between-group separation already shown in the main scenario simulations, but does not itself create an additional or different network-identification failure beyond what unmodeled, within-group context heterogeneity already produces. A context-adjusted estimator remains an effective correction whether the underlying context heterogeneity within a group arises from simple sampling variation, as in the main check, or from the model's own coupled feedback dynamics, as here.

\begin{figure}[htbp]
  \centering
  \includegraphics[width=\linewidth]{figs/Figure_S_network_feedback_check.pdf}
  \caption{\textbf{Supplementary feedback check.} Each simulated individual's slow context was generated by their own private, coupled fast/slow system, using the same calibrated dynamics as the main scenario simulations, with feedback either disabled ($b=0$) or enabled ($b=0.15$). \textbf{(A)} Realized between-person standard deviation of the final context value, $SD(P)$, by group and feedback condition. \textbf{(B)} Estimated global strength for the symptom-only and context-adjusted estimators; the dotted line shows the true global strength. \textbf{(C)} Mean absolute recovery error relative to $W^{\mathrm{true}}$ for both estimators.}
  \label{fig:network_feedback_check}
\end{figure}
